% explainer-source-hash: 378da11993b49d26f51eace8aae858845ac55c6ff180b36302677526a7ffcf01
% explainer-source-commit: b46124dd2204fce0b088f79ad7cd690afbbcf024
% explainer-generated: 2026-07-16
% Regenerate with /explain-all (see WIRING.md). Do not edit this stamp by hand:
% it is how the toolchain knows whether this document still matches the code.
\documentclass[11pt,a4paper]{article}
\input{preamble}
\begin{document}

\begin{titlepage}
\centering
\vspace*{3cm}
{\Huge\bfseries\color{inkblue} Morse Code Converter\par}
\vspace{0.8cm}
{\Large\color{accent} The Deep Dive\par}
\vspace{1.2cm}
{\large An exhaustive, beginner-safe walkthrough of every file,\\
every function, and every design decision in the project.\par}
\vspace{2cm}
\begin{minipage}{0.8\textwidth}
\centering\small\itshape
This document assumes no prior programming knowledge. Every technical term is
defined at the point where it first appears. Everything stated here was read out
of the source code or confirmed by running it; anything inferred is labelled as
inferred.
\end{minipage}
\vfill
\end{titlepage}

\tableofcontents
\newpage

% =====================================================================
\section{What this project is}

The Morse Code Converter is a small desktop application, written in Python, that
turns ordinary text (\code{SOS}) into International Morse Code
(\code{... --- ...}) and turns Morse code back into ordinary text. It
does this live, as you type, in a window with two tabs. It can also play the
resulting Morse code out loud as real beeps, while flashing each dot and dash on
screen in time with the sound.

The whole project is three Python files and roughly 600 lines of code including
comments. It has no third-party dependencies whatsoever: everything it does, it
does with what ships inside Python itself.

\begin{jargon}{Morse code}
A way of representing letters and numbers as sequences of two symbols: a short
signal (a \emph{dot}, written \code{.}) and a long signal (a \emph{dash}, written
\code{-}). It was designed in the 1830s for sending messages down a telegraph
wire, where the only thing you could control was whether current was flowing or
not. The letter S, for example, is three dots (\code{...}) and the letter O is
three dashes (\code{---}), which is why the distress call SOS is written
\code{... --- ...}. ``International''
Morse Code (also called ITU Morse) is the modern standardised version, as opposed
to the older American Morse.
\end{jargon}

\begin{jargon}{GUI (Graphical User Interface)}
The part of a program you see and click: windows, buttons, text boxes. The
opposite is a \emph{command-line interface} (CLI), where you type a command into a
terminal and read text back. This project has both, but the GUI is the main event.
\end{jargon}

\begin{jargon}{Module}
A single Python file containing code. One module can \emph{import} another, which
means ``load that file and let me use the things defined inside it''. Splitting a
program into modules is how you keep unrelated concerns in separate places.
\end{jargon}

\begin{keyidea}{Why the project is shaped the way it is}
The single organising decision is that the \textbf{conversion logic knows nothing
about the interface}, and the \textbf{audio logic knows nothing about the
interface either}. Two of the three files can be imported and used with no window
ever opening. Only the third file, \file{gui.py}, knows that a screen exists. This
is what makes the logic testable and the audio reusable, and it is the first thing
to point at when explaining the project.
\end{keyidea}

% =====================================================================
\section{The shape of the repository}

\begin{figure}[H]
\centering
\begin{forest} dirtree
[morse-code-converter
  [morse\_converter.py {\rmfamily\itshape \quad the alphabet and the two conversions (no UI, no audio)}]
  [morse\_audio.py {\rmfamily\itshape \quad timing, tone generation, playback backends (no UI)}]
  [gui.py {\rmfamily\itshape \quad the tkinter window; imports the other two}]
  [requirements.txt {\rmfamily\itshape \quad declares that there are no dependencies}]
  [README.md {\rmfamily\itshape \quad project documentation}]
  [LICENSE {\rmfamily\itshape \quad PolyForm Strict 1.0.0}]
  [.gitignore {\rmfamily\itshape \quad files git should not track}]
  [docs/
    [screenshots/
      [app.png {\rmfamily\itshape \quad a screenshot of the running window}]
    ]
    [explainers/
      [deep-dive.tex {\rmfamily\itshape \quad this document}]
      [brief.tex {\rmfamily\itshape \quad the shorter companion}]
    ]
  ]
]
\end{forest}
\caption{Every tracked file in the project. Three of them contain code.}
\end{figure}

The file-organisation rationale is visible directly in the import statements: the
arrows only ever point one way.

\begin{figure}[H]
\centering
\begin{tikzpicture}[node distance=10mm]
  \node[box] (gui) {\file{gui.py}\\[2pt] \scriptsize the tkinter window\\ \scriptsize tabs, boxes, canvas, buttons};
  \node[box, below left=16mm and 6mm of gui] (conv) {\file{morse\_converter.py}\\[2pt] \scriptsize \code{MORSE\_CODE} table\\ \scriptsize \code{text\_to\_morse}\\ \scriptsize \code{morse\_to\_text}};
  \node[box, below right=16mm and 6mm of gui] (audio) {\file{morse\_audio.py}\\[2pt] \scriptsize \code{build\_timing\_sequence}\\ \scriptsize tone generation (WAV)\\ \scriptsize backend detection};
  \draw[flow] (gui) -- node[lbl,left] {imports} (conv);
  \draw[flow] (gui) -- node[lbl,right] {imports} (audio);
  \node[below=13mm of conv, font=\scriptsize\itshape, text=black!65] {no UI, no audio, no I/O};
  \node[below=13mm of audio, font=\scriptsize\itshape, text=black!65] (anote) {no UI, standard library only};
  \node[extbox, below=7mm of anote] (os) {operating system\\[2pt] \scriptsize \code{aplay} / \code{paplay}\\ \scriptsize \code{afplay} / \code{ffplay}\\ \scriptsize \code{winsound}};
  \draw[dflow] (audio.east) -- ++(8mm,0) coordinate (bend) |- (os.east);
  \node[lbl] at ($(bend)!0.5!(bend|-os.east)$) {spawns};
\end{tikzpicture}
\caption{Module dependency graph. Nothing points upward: neither
\file{morse\_converter.py} nor \file{morse\_audio.py} imports \file{gui.py}, and
they do not import each other in normal use.}
\end{figure}

\begin{honest}{One import that crosses the line, in a corner}
\file{morse\_audio.py} does import \file{morse\_converter.py}, but only \emph{inside}
its \code{if \_\_name\_\_ == "\_\_main\_\_":} block, which runs solely when you execute
that file directly as a smoke test. So the clean separation holds for every real
code path, and the exception is deliberate and contained. It is worth knowing about
rather than claiming the modules are perfectly independent.
\end{honest}

\begin{jargon}{\code{if \_\_name\_\_ == "\_\_main\_\_":}}
A standard Python idiom. Python sets the variable \code{\_\_name\_\_} to the string
\code{"\_\_main\_\_"} when a file is run directly (\code{python3 gui.py}), and to the
file's own name when the file is merely imported by another file. So code guarded by
this line runs only in the first case. It is how one file can be both an importable
library and a runnable script.
\end{jargon}

% =====================================================================
\section{\file{morse\_converter.py}: the alphabet and the two conversions}

This is the smallest and most important file: 62 lines, no imports at all.

\subsection{\code{MORSE\_CODE}: the lookup table}

\begin{jargon}{Dictionary (in Python: \code{dict})}
A container that stores pairs, mapping a \emph{key} to a \emph{value}, and lets you
retrieve the value instantly by naming the key. Think of a physical dictionary:
you look up ``aardvark'' (the key) and get its definition (the value). Crucially,
the lookup does not get slower as the dictionary grows, because it uses a
\emph{hash table} internally: it computes a number from the key and jumps straight
to the right slot rather than scanning every entry.
\end{jargon}

The heart of the program is a hardcoded dictionary mapping each supported character
to its Morse pattern:

% literate= is load-bearing: listings + inputenc cannot read the two multi-byte
% characters in this table (inverted question mark and inverted exclamation mark)
% on their own; without the mapping below the build dies with
% "LaTeX Error: Invalid UTF-8 byte sequence". They are real entries in the real
% table, so the fix is to map them for the renderer, never to drop them.
\begin{lstlisting}[language=Python,
  literate={¿}{{\textquestiondown}}1 {¡}{{\textexclamdown}}1,
  caption={The complete table, reproduced exactly as it appears in \file{morse\_converter.py}. This is the whole of the project's domain knowledge: every other line of code exists to move data through this dictionary.}]
MORSE_CODE = {
    "A": ".-", "B": "-...", "C": "-.-.", "D": "-..", "E": ".",
    "F": "..-.", "G": "--.", "H": "....", "I": "..", "J": ".---",
    "K": "-.-", "L": ".-..", "M": "--", "N": "-.", "O": "---",
    "P": ".--.", "Q": "--.-", "R": ".-.", "S": "...", "T": "-",
    "U": "..-", "V": "...-", "W": ".--", "X": "-..-", "Y": "-.--",
    "Z": "--..",
    "0": "-----", "1": ".----", "2": "..---", "3": "...--", "4": "....-",
    "5": ".....", "6": "-....", "7": "--...", "8": "---..", "9": "----.",
    " ": "/",
    "!": "-.-.--", "?": "..--..", ".": ".-.-.-", ",": "--..--",
    "'": ".----.", '"': ".-..-.", "/": "-..-.", "(": "-.--.",
    ")": "-.--.-", "&": ".-...", ":": "---...", ";": "-.-.-.",
    "=": "-...-", "+": ".-.-.", "-": "-....-", "_": "..--.-",
    "$": "...-..-", "@": ".--.-.", "¿": "..-.-", "¡": "--...-",
}
\end{lstlisting}

The table holds \textbf{57 entries} (confirmed by loading the module and counting):
26 letters, 10 digits, the space character, and 20 punctuation marks including the
inverted Spanish \code{¿} and \code{¡}.

Two rows repay a second look. \code{"O": "---"} is three dashes, the longest of the
plain letters alongside \code{"0": "-----"}, which is five; and the hyphen character
itself is an entry, \code{"-": "-....-"}, so a dash can be spelled in dashes. Both
are worth noticing because they are the rows that a document about Morse code is
most tempted to quietly skip.

The most consequential row is \code{" ": "/"}. The space character is treated as
just another character in the alphabet, whose Morse ``pattern'' is a forward slash.
This one decision quietly does a lot of work, as we will see in
Section~\ref{sec:redundant}.

\subsection{\code{TEXT\_CODE}: the table, inverted}

\begin{lstlisting}[language=Python]
TEXT_CODE = {value: key for key, value in MORSE_CODE.items()}
\end{lstlisting}

\begin{jargon}{Dictionary comprehension}
A compact way of building a dictionary from another collection in one line. Read
\code{\{value: key for key, value in MORSE\_CODE.items()\}} as: ``go through every
key/value pair in \code{MORSE\_CODE}, and build a new dictionary in which each old
value becomes the new key and each old key becomes the new value''. The result is
the same table read backwards: \code{".-"} now maps to \code{"A"}.
\end{jargon}

This is a neat trick, and it means there is exactly one place to edit if a mapping
is wrong. But it carries a silent hazard.

\begin{keyidea}{The inversion hazard, and why this project escapes it}
Dictionary keys must be unique. If two different characters had the \emph{same}
Morse pattern, inverting the table would silently discard one of them: no error, no
warning, just a decoder that is quietly wrong for one character forever. Nothing in
the code checks for this.

I checked it directly rather than assuming. Loading the module and counting gives
\textbf{57 entries in \code{MORSE\_CODE} and 57 in \code{TEXT\_CODE}}, with zero
duplicate values. So the inversion is lossless today, and the code is correct. The
hazard is that it is correct \emph{by luck of the table's contents}, not by any
guard, and a future edit that introduced a duplicate would not be caught.
\end{keyidea}

\subsection{\code{text\_to\_morse}: encoding}

\begin{lstlisting}[language=Python]
def text_to_morse(text: str) -> str:
    """Convert a string of text to Morse code, characters separated by spaces."""
    codes = []
    for char in text.upper():
        if char not in MORSE_CODE:
            raise ValueError(f"No Morse mapping for character: {char!r}")
        codes.append(MORSE_CODE[char])
    return " ".join(codes)
\end{lstlisting}

Line by line:

\begin{itemize}
\item \code{text: str -> str} is a \emph{type hint}: a note to the reader (and to
      automated checkers) that this function takes a string and returns a string.
      Python does not enforce it at runtime.
\item \code{text.upper()} converts the whole input to capitals first. This is the
      entire implementation of the ``case-insensitive'' feature: the table only
      contains capitals, so lowercase input is normalised before it is ever looked
      up.
\item The \code{if char not in MORSE\_CODE} check is the deliberate part. Without
      it, \code{MORSE\_CODE[char]} on an unknown character would raise a
      \code{KeyError}, a low-level error naming only the character with no context.
      Instead the function raises a \code{ValueError} carrying a human-readable
      message. The \code{!r} inside the f-string means ``show this using
      \code{repr}'', so an invisible character prints as \code{'\textbackslash n'}
      rather than as an actual blank line.
\item \code{" ".join(codes)} glues the per-character patterns together with a
      single space between each.
\end{itemize}

\begin{jargon}{Raising an exception}
When a function cannot do its job, it can \emph{raise} an exception: it stops
immediately and hands an error object back up to whoever called it. The caller can
\emph{catch} it (with \code{try}/\code{except}) and react, or let it propagate and
crash the program. \code{ValueError} conventionally means ``you gave me a value I
cannot work with''.
\end{jargon}

So \code{text\_to\_morse("SOS")} walks S, O, S and produces \code{... --- ...}: the
canonical Morse message, three dots, three dashes, three dots. Add a second word and
the space rule shows itself: \code{text\_to\_morse("SOS SOS")} gives
\code{... --- ... / ... --- ...}. Note that the space became \code{/} through the
ordinary table lookup, and then got a space on each side from the \code{join}, which
is why every word boundary in the output is literally the three characters
\code{"\ /\ "}.

\subsection{\code{morse\_to\_text}: decoding}

\begin{lstlisting}[language=Python]
def morse_to_text(morse: str) -> str:
    words = morse.strip().split(" / ")
    decoded_words = []
    for word in words:
        letters = []
        for symbol in word.split():
            if symbol not in TEXT_CODE:
                raise ValueError(f"No text mapping for Morse symbol: {symbol!r}")
            letters.append(TEXT_CODE[symbol])
        decoded_words.append("".join(letters))
    return " ".join(decoded_words)
\end{lstlisting}

This is the mirror image, in two nested levels:

\begin{figure}[H]
\centering
\begin{tikzpicture}[node distance=8mm]
  \node[box] (in) {\code{"... --- ... / ... --- ..."}};
  \node[box, below=of in] (split1) {\code{.strip().split(" / ")}\\[2pt] \scriptsize split into words};
  \node[box, below=of split1] (words) {\code{["... --- ..."]}\quad\code{["... --- ..."]}};
  \node[box, below=of words] (split2) {\code{word.split()}\\[2pt] \scriptsize split into letter symbols};
  \node[box, below=of split2] (syms) {\code{"..."} \code{"---"} \code{"..."}};
  \node[box, below=of syms] (look) {look each up in \code{TEXT\_CODE}\\[2pt] \scriptsize unknown symbol $\Rightarrow$ \code{ValueError}};
  \node[box, below=of look] (out) {\code{"SOS"} + \code{"SOS"} joined with a space\\[2pt] \code{"SOS SOS"}};
  \draw[flow] (in) -- (split1);
  \draw[flow] (split1) -- (words);
  \draw[flow] (words) -- (split2);
  \draw[flow] (split2) -- (syms);
  \draw[flow] (syms) -- (look);
  \draw[flow] (look) -- (out);
\end{tikzpicture}
\caption{The decoding pipeline for \code{"SOS SOS"}: split on words, then on
letters, then look up each symbol. Both words decode through the same three
lookups, \code{"..."} $\rightarrow$ S, \code{"---"} $\rightarrow$ O,
\code{"..."} $\rightarrow$ S.}
\end{figure}

\code{.strip()} removes leading and trailing whitespace so a pasted string with a
stray newline still decodes. \code{word.split()} with no argument splits on any run
of whitespace, so extra spaces between letters are tolerated.

\subsection{The round-trip property, tested}

The docstring claims that \code{morse\_to\_text(text\_to\_morse(s)) == s.upper()}.
That is a strong claim, so I tested it rather than repeating it. Running the two
functions over letters, digits, punctuation, mixed case, leading and trailing
spaces, doubled internal spaces, an all-whitespace string, and a string containing
a literal forward slash (\code{"ABC/DEF"}, where \code{/} is itself a table entry
mapping to \code{-..-.}) gave a correct round trip in every case. The claim holds
for the inputs the encoder accepts.

\subsection{The redundant word-splitting layer}
\label{sec:redundant}

\begin{honest}{The outer \code{split(" / ")} does nothing the table was not already doing}
\code{morse\_to\_text} handles word gaps twice. It splits the input on \code{" / "},
decodes each word separately, and re-joins with spaces. But \code{"/"} is
\emph{also} an ordinary entry in \code{TEXT\_CODE}, mapping to the space character,
because \code{" ": "/"} is in \code{MORSE\_CODE}. So the simpler one-liner

\begin{lstlisting}[language=Python]
"".join(TEXT_CODE[s] for s in morse.strip().split())
\end{lstlisting}

produces identical output. I confirmed this empirically: across the same test
inputs above (including the awkward leading-space, trailing-space and
doubled-space cases) the naive version and the shipped version diverged on
\textbf{zero} inputs.

The README's ``Challenges'' section defends the split-on-\code{" / "} approach with
``splitting naively on \code{/} alone would have left stray spaces around each
word''. That reasoning is about splitting on the slash, but the alternative it
overlooks is not splitting on the slash at all, and simply \emph{decoding} it like
any other symbol. Two mechanisms now encode the same rule, and they have to agree
forever. This is not a bug and it does not misbehave; it is redundant complexity
with a rationale in the README that does not survive contact with the code.
\end{honest}

\subsection{The command-line fallback}

\begin{lstlisting}[language=Python]
if __name__ == "__main__":
    message = input("Enter a message to be converted into Morse code: ")
    print(text_to_morse(message))
\end{lstlisting}

Running \code{python3 morse\_converter.py} gives a one-shot prompt. Note it is
\textbf{one-directional}: the CLI can encode but not decode, even though
\code{morse\_to\_text} sits right above it in the same file. Note also that it does
not catch \code{ValueError}, so typing an unsupported character here produces a raw
Python traceback. Both are fine for a development convenience, and the README is
upfront that this is a leftover from the project's origin as a bare script.

% =====================================================================
\section{\file{morse\_audio.py}: turning dots into sound}

This is the most technically interesting file, and the one whose central claim
(``real audio with no dependencies'') is worth understanding properly.

\subsection{The timing model}

Morse code is not just a sequence of dots and dashes; it is a sequence of
\emph{durations}. The standard defines everything in terms of one base unit:

\begin{table}[H]
\centering
\begin{tabular}{lll}
\toprule
\textbf{Element} & \textbf{Duration} & \textbf{Constant in the code} \\
\midrule
dot & 1 unit of tone & \code{DOT\_UNITS = 1} \\
dash & 3 units of tone & \code{DASH\_UNITS = 3} \\
gap between symbols of one letter & 1 unit of silence & \code{INTRA\_SYMBOL\_GAP\_UNITS = 1} \\
gap between letters & 3 units of silence & \code{INTER\_LETTER\_GAP\_UNITS = 3} \\
gap between words & 7 units of silence & \code{INTER\_WORD\_GAP\_UNITS = 7} \\
\bottomrule
\end{tabular}
\caption{The standard Morse timing ratios, each given a named constant at the top
of \file{morse\_audio.py}. One unit is set to 80~ms (\code{UNIT\_MS = 80}), which
fixes the sending speed.}
\end{table}

\begin{keyidea}{Why the three different gaps matter}
Silence carries information in Morse code. Without three distinct gap lengths, a
listener cannot tell where one letter ends and the next begins, and
\code{.... ..} (``HI'') would be indistinguishable from a single six-dot blob. The
gaps are not padding, they are punctuation.
\end{keyidea}

\subsection{\code{build\_timing\_sequence}: the single source of truth}

This function converts a Morse string into an explicit list of
\code{(kind, units)} pairs, where \code{kind} is \code{"dot"}, \code{"dash"} or
\code{"gap"}.

\begin{pseudocode}{build\_timing\_sequence(morse)}
\Step{events := empty list}
\Step{words := morse.strip() split on " / ", dropping empty pieces}
\Step{for each word, with index w:}
\Indent{letters := word split on " ", dropping empty pieces}
\Indent{for each letter, with index l:}
\Indent{\hspace*{2em}symbols := the characters of letter that are "." or "-"}
\Indent{\hspace*{2em}for each symbol, with index s:}
\Indent{\hspace*{4em}emit ("dot",1) if ".", else ("dash",3)}
\Indent{\hspace*{4em}if s is not the last symbol: emit ("gap",1)}
\Indent{\hspace*{2em}if l is not the last letter: emit ("gap",3)}
\Indent{if w is not the last word: emit ("gap",7)}
\Step{return events}
\end{pseudocode}

The three ``if not the last'' checks are the whole trick. Gaps are emitted
\emph{between} elements, never after the final one, so a message never ends with
trailing silence and gaps never double up at a boundary.

For \code{"SOS"} this produces 17 events totalling \textbf{27 units}. Repeating the
call as \code{"SOS SOS"}, which is how a distress signal is actually sent, produces
35 events totalling \textbf{61 units} and adds the one thing a single word cannot
show: the 7-unit word gap. Here is the real timeline it generates, drawn to scale:

\begin{figure}[H]
\centering
\begin{tikzpicture}[x=1.2mm, y=1mm]
  % "SOS SOS" = ... --- ... / ... --- ...
  % S = 3 dots + two 1u gaps = 5u; O = 3 dashes + two 1u gaps = 11u
  % word 1: S 0-5, gap3, O 8-19, gap3, S 22-27 ; word gap 7 ; word 2: S 34-39, gap3, O 42-53, gap3, S 56-61
  % ---- word 1: S
  \fill[accent] (0,0) rectangle (1,6);
  \fill[accent] (2,0) rectangle (3,6);
  \fill[accent] (4,0) rectangle (5,6);
  \draw[decorate, decoration={brace, amplitude=3pt}, draw=inkblue] (0,7.5) -- (5,7.5) node[midway, above, font=\scriptsize] {S};
  % ---- word 1: O  (three dashes)
  \fill[inkblue] (8,0) rectangle (11,6);
  \fill[inkblue] (12,0) rectangle (15,6);
  \fill[inkblue] (16,0) rectangle (19,6);
  \draw[decorate, decoration={brace, amplitude=3pt}, draw=inkblue] (8,7.5) -- (19,7.5) node[midway, above, font=\scriptsize] {O};
  % ---- word 1: S
  \fill[accent] (22,0) rectangle (23,6);
  \fill[accent] (24,0) rectangle (25,6);
  \fill[accent] (26,0) rectangle (27,6);
  \draw[decorate, decoration={brace, amplitude=3pt}, draw=inkblue] (22,7.5) -- (27,7.5) node[midway, above, font=\scriptsize] {S};
  % ---- word gap, 7 units
  \draw[decorate, decoration={brace, amplitude=3pt, mirror}, draw=warnred] (27,-1.5) -- (34,-1.5) node[midway, below, font=\scriptsize, text=warnred] {word gap, 7u};
  % ---- word 2: S
  \fill[accent] (34,0) rectangle (35,6);
  \fill[accent] (36,0) rectangle (37,6);
  \fill[accent] (38,0) rectangle (39,6);
  \draw[decorate, decoration={brace, amplitude=3pt}, draw=inkblue] (34,7.5) -- (39,7.5) node[midway, above, font=\scriptsize] {S};
  % ---- word 2: O
  \fill[inkblue] (42,0) rectangle (45,6);
  \fill[inkblue] (46,0) rectangle (49,6);
  \fill[inkblue] (50,0) rectangle (53,6);
  \draw[decorate, decoration={brace, amplitude=3pt}, draw=inkblue] (42,7.5) -- (53,7.5) node[midway, above, font=\scriptsize] {O};
  % ---- word 2: S
  \fill[accent] (56,0) rectangle (57,6);
  \fill[accent] (58,0) rectangle (59,6);
  \fill[accent] (60,0) rectangle (61,6);
  \draw[decorate, decoration={brace, amplitude=3pt}, draw=inkblue] (56,7.5) -- (61,7.5) node[midway, above, font=\scriptsize] {S};
  % axis
  \draw[draw=linegrey] (0,-0.3) -- (61,-0.3);
  % legend
  \fill[accent] (0,-10) rectangle (1.6,-7.8);
  \node[font=\scriptsize, anchor=west] at (3,-8.9) {dot (1u)};
  \fill[inkblue] (16,-10) rectangle (17.6,-7.8);
  \node[font=\scriptsize, anchor=west] at (19,-8.9) {dash (3u)};
  \draw[draw=linegrey] (34,-10) rectangle (35.6,-7.8);
  \node[font=\scriptsize, anchor=west] at (37,-8.9) {silence (gap)};
  \node[font=\scriptsize, text=black!65, anchor=west] at (0,-13.5) {61 units $\times$ 80\,ms = 4880\,ms};
\end{tikzpicture}
\caption{The real event timeline \code{build\_timing\_sequence} emits for
\code{"SOS SOS"}, drawn to scale. Each S is three 1-unit dots and each O is three
3-unit dashes, separated by 1-unit gaps inside the letter and 3-unit gaps between
letters, with the single 7-unit word gap in the middle. The dash being visibly three
times the dot is the ratio the project originally got wrong. This diagram was drawn
from the function's actual output, not from the specification.}
\end{figure}

\begin{keyidea}{One list drives both the eyes and the ears}
This event list is the project's best design decision. \file{gui.py} calls
\code{build\_timing\_sequence} \textbf{once} and uses the result for both the
on-screen flashing and the tone playback. Because both read the same list, the
sound and the picture cannot drift apart. The obvious alternative, hand-tuning a
visual delay and an audio delay separately, is exactly what the project used to do,
and the README records that it produced a dot-to-dash ratio of roughly 1:1.9
instead of the correct 1:3.
\end{keyidea}

\begin{honest}{Silent tolerance of garbage, inconsistent with the decoder}
The line \code{symbols = [c for c in letter if c in ".-"]} keeps only dots and
dashes and \textbf{silently discards everything else}. Feeding
\code{build\_timing\_sequence(".- XYZ -...")} does not raise; it drops \code{XYZ}
entirely but still emits the inter-letter gap around it, producing a timeline with
two adjacent 3-unit gaps and no symbol between them (I confirmed this by calling
the function directly).

So the two halves of the program disagree about invalid input:
\code{morse\_to\_text} raises a \code{ValueError} naming the bad symbol, while
\code{build\_timing\_sequence} shrugs and carries on. In the GUI this is masked,
because invalid Morse never reaches playback via the normal path, but it is a real
inconsistency in the library's contract and worth owning rather than explaining
away.
\end{honest}

\subsection{Generating a tone from nothing}

\begin{jargon}{Digital audio, PCM, and sample rate}
A sound is a wave of air pressure. A computer stores it by measuring that pressure
many thousands of times per second and writing down each measurement as a number.
Each measurement is a \emph{sample}; the number of measurements per second is the
\emph{sample rate} (here \code{SAMPLE\_RATE = 44100}, or 44{,}100 per second, the
same rate used by audio CDs). This raw list of numbers is called \emph{PCM}
(Pulse-Code Modulation). ``16-bit'' means each number is stored in 16 bits, giving
a range of $-32768$ to $+32767$.
\end{jargon}

A pure musical tone is a \emph{sine wave}: a value that swings smoothly up and down
at a fixed rate. The number of swings per second is the \emph{frequency}, measured
in hertz (Hz), and it is what we perceive as pitch. This project uses 700~Hz
(\code{TONE\_FREQUENCY\_HZ = 700}), a typical Morse practice pitch.

\begin{lstlisting}[language=Python, caption={The entire tone generator.}]
def _sine_pcm_frames(duration_ms, frequency=TONE_FREQUENCY_HZ,
                      sample_rate=SAMPLE_RATE, volume=VOLUME):
    """Raw 16-bit little-endian PCM bytes for a pure sine tone."""
    n_samples = int(sample_rate * duration_ms / 1000)
    amplitude = int(volume * 32767)
    samples = [
        int(amplitude * math.sin(2 * math.pi * frequency * (i / sample_rate)))
        for i in range(n_samples)
    ]
    return struct.pack("<%dh" % len(samples), *samples)
\end{lstlisting}

This is the crux of the ``no dependencies'' claim, so it is worth reading closely:

\begin{itemize}
% \sloppy is local to this list: these items are dense with unbreakable \code spans,
% so without it the first line of the math.sin item runs past the right margin.
\sloppy
\item \code{n\_samples} is how many measurements this tone needs: 80~ms at 44100
      samples per second is 3528 samples.
\item \code{amplitude} scales the wave's height. \code{VOLUME = 0.5} means the tone
      peaks at half the maximum a 16-bit sample can hold, leaving headroom.
\item \code{math.sin(2 * math.pi * frequency * (i / sample\_rate))} is the sine
      wave itself. \code{i / sample\_rate} converts the sample number into a time in
      seconds; multiplying by the frequency and by $2\pi$ (one full cycle in
      radians) gives the point on the wave at that moment.
\item \code{struct.pack("<\%dh" \% len(samples), *samples)} converts the Python list
      of numbers into raw bytes. \code{<} means little-endian (least significant
      byte first, the byte order WAV files use), and \code{h} means ``signed
      2-byte integer''. The \code{\%d} is filled in with the count, so the format
      string becomes e.g. \code{"<3528h"}.
\end{itemize}

\begin{jargon}{The \code{wave} module and the WAV format}
A WAV file is essentially raw PCM samples with a small header on the front
declaring the sample rate, how many channels there are (1 = mono, 2 = stereo), and
how many bytes each sample uses. Python's standard library includes \code{wave},
which writes that header for you. This is why the project needs no audio library:
the file format is simple enough that the standard library already covers it.
\end{jargon}

\code{write\_tone\_wav} is the thin wrapper that puts those two together: generate
the frames, open a WAV file, declare mono (\code{setnchannels(1)}), 2 bytes per
sample (\code{setsampwidth(2)}), the sample rate, and write.

\subsection{The tone cache}

\begin{lstlisting}[language=Python]
_tmp_dir = None

def _tone_dir():
    global _tmp_dir
    if _tmp_dir is None:
        _tmp_dir = tempfile.mkdtemp(prefix="morse_tones_")
        atexit.register(shutil.rmtree, _tmp_dir, ignore_errors=True)
    return _tmp_dir
\end{lstlisting}

The pattern here is \emph{lazy initialisation}: the temporary directory is created
the first time it is needed, not at import time, and only once. \code{atexit.register}
schedules the directory's deletion for when Python exits, so the program cleans up
after itself without anyone having to remember to call a cleanup function.

\code{build\_symbol\_tones} then writes exactly two files into that directory,
\file{dot.wav} (80~ms of tone) and \file{dash.wav} (240~ms of tone), once per app
run. Every dot in every message replays the same file. This is a sensible trade:
two small files generated once, rather than a fresh sine wave computed for every
symbol.

\subsection{Finding something that can actually make a sound}

Generating a WAV is the easy half. \emph{Playing} it is where portability bites,
because Python's standard library has no cross-platform way to send audio to a
speaker. The project's answer is to ask the machine what it has, at runtime.

\begin{figure}[H]
\centering
\begin{tikzpicture}[node distance=7mm]
  \node[box] (start) {\code{detect\_playback\_backend()}};
  \node[dec, below=of start] (win) {Windows?};
  \node[box, right=16mm of win] (ws) {import \code{winsound}};
  \node[dec, below=8mm of ws] (wsok) {imported?};
  \node[box, right=12mm of wsok, fill=okgreen!12] (retws) {return\\ \code{"winsound"}};
  \node[dec, below=14mm of win] (aplay) {\code{which("aplay")}};
  \node[dec, below=of aplay] (paplay) {\code{which("paplay")}};
  \node[dec, below=of paplay] (afplay) {\code{which("afplay")}};
  \node[dec, below=of afplay] (ffplay) {\code{which("ffplay")}};
  \node[box, below=of ffplay, fill=warnred!10] (none) {return \code{None}\\[2pt] \scriptsize no audio on this machine};
  \node[box, right=30mm of paplay, fill=okgreen!12] (found) {return that\\ player's name};

  \draw[flow] (start) -- (win);
  \draw[flow] (win) -- node[lbl,above] {yes} (ws);
  \draw[flow] (ws) -- (wsok);
  \draw[flow] (wsok) -- node[lbl,above] {yes} (retws);
  \draw[flow] (win) -- node[lbl,left] {no} (aplay);
  \draw[flow] (wsok) -| node[lbl,pos=0.25,below] {ImportError} (aplay);
  \draw[flow] (aplay) -- node[lbl,left] {not found} (paplay);
  \draw[flow] (paplay) -- node[lbl,left] {not found} (afplay);
  \draw[flow] (afplay) -- node[lbl,left] {not found} (ffplay);
  \draw[flow] (ffplay) -- node[lbl,left] {not found} (none);
  \draw[flow] (aplay) -| node[lbl,pos=0.25,above] {found} (found);
  \draw[flow] (paplay) -- node[lbl,above] {found} (found);
  \draw[flow] (afplay) -| node[lbl,pos=0.25,below] {found} (found);
  \draw[flow] (ffplay) -| node[lbl,pos=0.25,below] {found} (found);
\end{tikzpicture}
\caption{Backend detection. Nothing is assumed about the host: every candidate is
checked for real, and \code{None} is a legitimate, handled answer.}
\end{figure}

\begin{jargon}{\code{shutil.which}}
A standard-library function that answers ``if I typed this command name in a
terminal, would anything run, and where is it?'' It searches the directories listed
in the \code{PATH} environment variable and returns the full path, or \code{None}.
It is the programmatic version of the Unix \code{which} command.
\end{jargon}

\begin{keyidea}{Detecting rather than assuming}
The alternative approach, ``I am on Linux, therefore \code{aplay} exists'', is
wrong often enough to matter: a minimal container may have neither ALSA nor
PulseAudio installed. By calling \code{shutil.which} and treating \code{None} as a
normal outcome, the program can tell the user the truth (``No audio output
available on this system, showing visual playback only'') instead of crashing or,
worse, silently doing nothing and looking broken.
\end{keyidea}

On this machine, I ran \code{detect\_playback\_backend()} and it returned
\code{'aplay'}, matching what the README claims.

\subsection{\code{play\_wav\_async}: fire and forget}

% sloppypar: the long argument lists are one unbreakable \code span each, and without
% it the second line of this paragraph runs past the right margin.
\begin{sloppypar}
Each backend needs a different command line, so the function maps the backend name
to its arguments (\code{["aplay", "-q", path]}, \code{["ffplay", "-nodisp",
"-autoexit", "-loglevel", "quiet", path]}, and so on), then launches it with
\code{subprocess.Popen}.
\end{sloppypar}

\begin{jargon}{Process, \code{subprocess.Popen}, and blocking}
A \emph{process} is a running program. \code{subprocess.Popen} starts another
program as a separate process and \textbf{returns immediately}, without waiting for
it to finish. Its cousins \code{subprocess.run} and \code{.wait()} instead
\emph{block}: they freeze the caller until the other program exits.
\end{jargon}

\begin{keyidea}{Why \code{Popen} and not \code{run} is load-bearing}
A GUI is a single loop constantly redrawing the window and reacting to input. If
that loop blocks for 240~ms waiting for a dash's tone process to exit, the window
freezes for 240~ms: the flash cannot turn off on time, and the picture falls behind
the sound. Using \code{Popen} keeps tkinter's own timer in charge of the whole
timeline, with audio fired off alongside it. The README identifies this correctly.
\end{keyidea}

\code{stdout} and \code{stderr} are redirected to \code{DEVNULL} (a black hole) so
the players cannot spray output into the terminal.

\begin{honest}{Fire and forget means never knowing}
The consequences of this design are real and the README states them: each symbol
pays the cost of spawning a whole operating-system process, and the program never
learns whether a tone actually played. The processes are never collected either
(the \code{Popen} object is returned and then discarded by the GUI), so on Unix
each finished player leaves a \emph{zombie} entry in the process table until the
Python process exits. For a program that plays a few dozen short tones this is
harmless, and the alternative (an audio device API) means a third-party dependency
the project deliberately refuses. It is a trade-off, honestly made, not an oversight.
\end{honest}

\subsection{\code{render\_morse\_to\_wav}: the offline renderer}

This function renders an \emph{entire} message, tones and silences alike, into one
WAV file, returning the event list it used.

\begin{honest}{Unused by the application}
Nothing in \file{gui.py} calls \code{render\_morse\_to\_wav}. Its only caller is the
\code{\_\_main\_\_} smoke test at the bottom of its own file. By a strict definition
it is dead code in the shipped application. The README is upfront that it exists for
offline rendering and was the tool used to verify timing during development, which is
a legitimate reason to keep it, and it earned its place again while I was checking
this project (see below). But it should be described as a development instrument,
not as an application feature.
\end{honest}

I used it to check the timing claim independently. Rendering \code{"SOS"} and
inspecting the resulting file gives \textbf{95{,}256 frames at 44{,}100~Hz, exactly
2160~ms}, matching the 27 units $\times$ 80~ms computed by hand from the
dot/dash/gap rules. Rendering \code{"SOS SOS"} gives \textbf{215{,}208 frames,
exactly 4880~ms}, matching 61 units $\times$ 80~ms, so the 7-unit word gap is
accounted for to the sample. The timing claim is confirmed, not merely asserted.

\begin{honest}{No fade on the tone edges (inferred, not observed)}
Each tone begins and ends abruptly at whatever point of the sine wave it happens to
reach; there is no short fade-in or fade-out (an \emph{envelope}). In audio work,
a hard cut mid-wave typically produces an audible click. I am flagging this as
\textbf{inferred from the code, not confirmed by listening}: I read the generator
and no envelope is applied, but I did not verify by ear whether a click is actually
perceptible at this volume and pitch. It is not a correctness problem in any case.
\end{honest}

% =====================================================================
\section{\file{gui.py}: the window}

\begin{jargon}{tkinter, widgets, and the event loop}
\emph{tkinter} is Python's built-in GUI toolkit, a wrapper around the older Tk
library. It ships with Python (though on Linux it sometimes needs a separate OS
package such as \code{python3-tk}). A \emph{widget} is any interface element: a
window, button, text box, canvas. \emph{ttk} is a newer set of themed widgets that
look more native. The \emph{event loop} (\code{root.mainloop()}) is an endless loop
that waits for things to happen (a key press, a click, a timer expiring) and calls
the function you registered for that event.
\end{jargon}

\begin{jargon}{Callback}
A function you hand to someone else so they can call it later, when something
happens. \code{command=self.\_start\_playback} on a button does not call
\code{\_start\_playback}; it stores it, to be called when the button is clicked.
\end{jargon}

\subsection{Construction order}

\code{MorseApp.\_\_init\_\_} does four things in a deliberate order:

\begin{enumerate}
\item Configure the window: title, dark background, 760$\times$560 starting size, and
      a \code{minsize} of 620$\times$480 so the layout cannot be crushed.
\item Set \code{\_playback\_job = None}, the handle used to cancel a scheduled
      animation step.
\item \textbf{Detect audio before building anything visible.} The result decides what
      the playback panel will say, so it has to be known first. If a backend exists,
      the two tone files are generated now, at startup, rather than during playback
      when the delay would be audible.
\item Build the styles, then the layout.
\end{enumerate}

\subsection{Styling}

\begin{sloppypar}
\code{\_build\_style} sets up a hand-rolled dark theme: a set of named colour
constants at the top of the file (\code{BG}, \code{PANEL\_BG}, \code{ACCENT} and so
on) applied through \code{ttk.Style}. It first tries \code{style.theme\_use("clam")}
inside a \code{try}/\code{except tk.TclError}, because \code{clam} is the theme that
actually honours custom colours, and the \code{except} means the app still opens on
a system where that theme is missing rather than dying at startup. This is a small
and well-judged piece of defensive code.
\end{sloppypar}

\subsection{The tab factory}

Rather than writing two nearly identical tabs, \code{\_build\_conversion\_tab} is
called twice with different arguments and returns a dictionary describing the tab it
just built:

\begin{lstlisting}[language=Python]
tab_state = {
    "frame": frame, "input_box": input_box, "output_box": output_box,
    "error_label": error_label, "convert_fn": convert_fn,
    "playback_source": playback_source,
}
\end{lstlisting}

The key idea is \code{convert\_fn}: the tab does not know \emph{which} direction it
converts. It just calls the function it was handed. The ``Text to Morse'' tab gets
one function, the ``Morse to Text'' tab gets the other, and the identical machinery
serves both.

\begin{jargon}{Passing a function as a value}
In Python, functions are ordinary values: they can be stored in a dictionary and
called later, exactly like a number or a string. This is what lets one piece of code
serve two directions.
\end{jargon}

\begin{honest}{Two wrapper functions that wrap nothing}
\begin{lstlisting}[language=Python]
def _safe_text_to_morse(self, raw: str) -> str:
    return text_to_morse(raw)

def _safe_morse_to_text(self, raw: str) -> str:
    return morse_to_text(raw)
\end{lstlisting}
These are pure pass-throughs. They add no behaviour at all, and the \code{\_safe\_}
prefix promises a safety they do not provide: neither catches anything, and the
actual error handling lives in \code{\_on\_key\_release}. A reader scanning the file
would reasonably assume these are where input is made safe, and be wrong. They could
be deleted and replaced with the imported functions directly. This is a small piece
of misleading dead abstraction.
\end{honest}

\subsection{Live conversion on every keystroke}

Each input box binds \code{<KeyRelease>} to \code{\_on\_key\_release}: after every
key comes back up, the whole box is re-read and re-converted. There is no submit
button because there is no submit.

\begin{figure}[H]
\centering
\begin{tikzpicture}[node distance=8mm]
  \node[box] (key) {user releases a key};
  \node[box, below=of key] (read) {read the whole input box\\ \code{get("1.0", "end-1c")}};
  \node[box, below=of read] (enable) {output box: \code{state="normal"}\\[2pt] \scriptsize temporarily writable};
  \node[dec, below=10mm of enable] (blank) {input\\ blank?};
  \node[box, right=18mm of blank] (clear) {clear output,\\ clear error,\\ re-disable, return};
  \node[box, below=12mm of blank] (try) {call \code{convert\_fn(raw)}};
  \node[dec, below=10mm of try] (ok) {raised\\ \code{ValueError}?};
  \node[box, left=16mm of ok, fill=warnred!10] (err) {result := \code{""}\\ error label := message};
  \node[box, below=12mm of ok] (write) {write result into output box,\\ clear error label};
  \node[box, below=of write] (disable) {output box: \code{state="disabled"}\\[2pt] \scriptsize read-only again};
  \draw[flow] (key) -- (read);
  \draw[flow] (read) -- (enable);
  \draw[flow] (enable) -- (blank);
  \draw[flow] (blank) -- node[lbl,above] {yes} (clear);
  \draw[flow] (blank) -- node[lbl,left] {no} (try);
  \draw[flow] (try) -- (ok);
  \draw[flow] (ok) -- node[lbl,above] {yes} (err);
  \draw[flow] (ok) -- node[lbl,left] {no} (write);
  \draw[flow] (err) |- (write);
  \draw[flow] (write) -- (disable);
\end{tikzpicture}
\caption{What happens on every single keystroke.}
\end{figure}

Two details are worth naming:

\begin{itemize}
\item \textbf{The disabled-output dance.} The output box is kept
      \code{state="disabled"} so the user cannot type into it. But a disabled tkinter
      \code{Text} rejects programmatic writes too, so the code must flip it to
      \code{"normal"}, write, and flip it back. That is why \code{state} is toggled
      three times in one short function.
\item \textbf{\code{"1.0"} and \code{"end-1c"}.} tkinter addresses text as
      ``line.column'', 1-indexed lines and 0-indexed columns, so \code{"1.0"} is the
      very start. \code{"end-1c"} means ``the end, minus one character'', because
      tkinter always appends an invisible trailing newline that you almost never want.
\item \textbf{Catching \code{ValueError} is the whole point.} Because conversion runs
      mid-word, invalid intermediate states are normal, not exceptional. The
      \code{except ValueError} turns what would be a crash into a small inline
      message.
\end{itemize}

\begin{honest}{A multi-line box that cannot accept multi-line input}
Both input boxes are \code{tk.Text} widgets with \code{height=6}: they look and
behave like multi-line editors, and pressing Enter inserts a newline. But
\code{"\textbackslash n"} is not in \code{MORSE\_CODE}, so the moment a user presses
Enter the conversion fails with \code{No Morse mapping for character:
'\textbackslash n'} and the output box goes blank. I confirmed this by calling
\code{text\_to\_morse("A\textbackslash n")} directly: it raises. Tab characters
behave the same way.

The failure is handled gracefully (an inline message, not a crash), so this is a
usability wrinkle rather than a defect: the widget invites input the alphabet cannot
represent. Mapping newline to a word separator, or stripping it before conversion,
would close the gap.
\end{honest}

\subsection{The playback engine}

Playback is the most intricate part of the GUI, and it is built as a small state
machine driven by tkinter's timer.

\begin{jargon}{\code{root.after(ms, fn)} and cooperative animation}
\code{after} asks tkinter to call \code{fn} once, roughly \code{ms} milliseconds
later, and returns a job handle that \code{after\_cancel} can use to call it off.
Nothing sleeps and nothing blocks: the event loop keeps running and simply fires the
callback when the time comes. Chaining \code{after} calls (each callback scheduling
the next) is the standard way to animate in tkinter without freezing the window.
\end{jargon}

\code{\_start\_playback} runs when Play is clicked:

\begin{enumerate}
\item Cancel any animation already in flight (\code{after\_cancel}), so a second
      click restarts cleanly rather than running two animations at once.
\item Ask \code{\_current\_tab\_state()} which Morse string to play.
\item Clear the canvas; if there is no Morse, stop here.
\item Call \code{build\_timing\_sequence(morse)} to get the event list.
\item \code{\_render\_events} lays out one shape per symbol on the canvas, dark
      (\code{SYMBOL\_OFF}), and pairs each with its real millisecond duration.
\item Set \code{\_play\_index = 0} and call \code{\_advance\_playback}.
\end{enumerate}

\code{\_render\_events} draws a small oval for each dot (14 pixels wide) and a
longer rectangle for each dash (34 pixels), so the shapes themselves encode the
distinction. Gaps get no shape, but they do advance the horizontal cursor by
\code{6 + units * 4} pixels, so a word gap opens a visibly wider space than a letter
gap. The visual spacing is proportional to the real timing but not on the same scale
as the shapes, which is a presentational choice, not a timing one.

\begin{figure}[H]
\centering
\begin{tikzpicture}[node distance=9mm]
  \node[box] (idle) {\textbf{idle}\\[2pt] \scriptsize \code{\_playback\_job = None}};
  \node[dec, below=10mm of idle] (more) {more\\ events?};
  \node[dec, below=12mm of more] (isgap) {kind ==\\ "gap"?};
  \node[box, right=22mm of isgap] (gapwait) {\code{after(duration)}\\[2pt] \scriptsize do nothing, just wait};
  \node[box, below=12mm of isgap] (on) {light the shape\\ (\code{DOT\_ON} / \code{DASH\_ON})\\[2pt] \scriptsize + \code{play\_wav\_async} if audio};
  \node[box, below=of on] (wait) {\code{after(duration, turn\_off)}};
  \node[box, below=of wait] (off) {turn the shape dark again\\ \code{\_play\_index += 1}};
  \draw[flow] (idle) -- node[lbl,left] {Play clicked} (more);
  \draw[flow] (more) -- node[lbl,left] {yes} (isgap);
  \draw[flow] (more) -| node[lbl,pos=0.25,above] {no: finished} ($(idle.east)+(3mm,0)$) -- (idle.east);
  \draw[flow] (isgap) -- node[lbl,above] {yes} (gapwait);
  \draw[flow] (isgap) -- node[lbl,left] {no} (on);
  \draw[flow] (on) -- (wait);
  \draw[flow] (wait) -- (off);
  \draw[flow] (off) -| node[lbl,pos=0.25,below] {recurse} ($(more.west)-(20mm,0)$) -- (more.west);
  \draw[flow] (gapwait) to[out=0,in=0,looseness=1.6] node[lbl,right] {\code{\_play\_index += 1}} (more.east);
\end{tikzpicture}
\caption{The playback state machine in \code{\_advance\_playback}. Each step
schedules the next one through \code{after}; the loop is a chain of timer callbacks,
never a blocking loop.}
\end{figure}

For a dot or dash, the sequence within one step is: light the shape, fire the tone
asynchronously, and schedule \code{turn\_off} for \code{duration\_ms} later.
\code{turn\_off} darkens the shape, increments the index, and calls
\code{\_advance\_playback} again. The audio and the light therefore start on the
same line of code and are extinguished by the same timer.

The \code{try}/\code{except OSError} around \code{play\_wav\_async} is deliberate:
if spawning the player fails partway through a message, the visual playback
continues uninterrupted. The code comment states the reasoning plainly, that the
flash is an honest representation on its own.

\begin{honest}{\code{\_current\_tab\_state} is misnamed and does not consider the current tab}
Despite the name, this function returns \textbf{a string}, not a tab state
dictionary, and it never asks the notebook which tab is actually selected. It checks
the Text-to-Morse tab's output box first and only falls back to the Morse-to-Text
tab's input box if the first is empty.

The consequence: type text in the first tab, switch to the second, type Morse there,
and press Play. It plays the \emph{first} tab's content, because that box is still
populated. The README's final ``Challenges'' bullet acknowledges the fixed priority
honestly and explains that the architecture diagram was redrawn to reflect it, which
is to the author's credit. But the priority itself is still surprising behaviour, and
the function name actively misdirects a reader about both its return type and its
logic.
\end{honest}

\begin{honest}{Attributes created outside \code{\_\_init\_\_}}
\code{\_playback\_events} and \code{\_play\_index} are first assigned inside
\code{\_start\_playback}, not in \code{\_\_init\_\_}. If \code{\_advance\_playback}
were ever called before a first Play click, it would fail with
\code{AttributeError}. No such path exists today (only \code{\_start\_playback}
begins the chain), so this is a latent fragility rather than a live bug, and
initialising both in \code{\_\_init\_\_} would cost one line.
\end{honest}

\subsection{The clipboard}

\code{\_copy\_to\_clipboard} uses \code{clipboard\_clear()} and
\code{clipboard\_append()}, both built into tkinter. This is a genuinely nice detail:
copying to the system clipboard is the classic reason a small Python project reaches
for a third-party package, and it turns out not to be necessary.

\begin{honest}{A known tkinter clipboard caveat (inferred, not tested here)}
On X11-based Linux systems, clipboard contents are owned by the application that set
them, and are conventionally lost when that application exits unless a clipboard
manager is running. So text copied from this app may not survive closing the window.
I am flagging this as \textbf{inferred from how X11 clipboards are documented to
work, not verified against this application}, and it affects nothing while the app is
open.
\end{honest}

% =====================================================================
\section{End to end: one keystroke and one click}

\begin{figure}[H]
\centering
\begin{tikzpicture}[node distance=11mm]
  \node[extbox] (user) {user};
  \node[box, below=of user] (input) {\code{tk.Text} input box};
  \node[box, below=of input] (bind) {\code{<KeyRelease>} $\rightarrow$ \code{\_on\_key\_release}};
  \node[box, below=of bind] (conv) {\code{text\_to\_morse} \\ \scriptsize (\file{morse\_converter.py})};
  \node[box, below=of conv] (out) {output \code{tk.Text} (read-only)};

  \node[box, right=34mm of bind] (play) {Play button $\rightarrow$ \code{\_start\_playback}};
  \node[box, below=of play] (seq) {\code{build\_timing\_sequence}\\ \scriptsize (\file{morse\_audio.py})};
  \node[box, below=of seq] (fork) {one event list};
  \node[box, below left=10mm and -6mm of fork] (canvas) {canvas flash\\ \scriptsize \code{after()} timer};
  \node[box, below right=10mm and -6mm of fork] (tone) {\code{play\_wav\_async}};
  \node[extbox, below=of tone] (spk) {system player\\ $\rightarrow$ speaker};

  \draw[flow] (user) -- node[lbl,left] {types} (input);
  \draw[flow] (input) -- (bind);
  \draw[flow] (bind) -- (conv);
  \draw[flow] (conv) -- node[lbl,left] {Morse string} (out);
  \draw[flow] (user.east) to[out=0,in=90] node[lbl,above] {clicks} (play.north);
  \draw[flow] (out.east) to[out=0,in=180] node[lbl,below,pos=0.4] {\code{\_current\_tab\_state()}\\ reads the box} (fork.west);
  \draw[flow] (play) -- (seq);
  \draw[flow] (seq) -- (fork);
  \draw[flow] (fork) -- (canvas);
  \draw[flow] (fork) -- (tone);
  \draw[dflow] (tone) -- node[lbl,right] {\code{Popen}} (spk);
\end{tikzpicture}
\caption{The complete data flow. The important feature is the fork at the bottom:
one event list feeds both the eyes and the ears, so they cannot disagree.}
\end{figure}

% =====================================================================
\section{Dependencies, configuration and licence}

\file{requirements.txt} contains one line, a comment stating there are no
third-party dependencies. The file exists so the answer to ``what does this need
installed?'' is explicit rather than absent.

Every import in the project is standard library:

\begin{table}[H]
\centering
\begin{tabular}{lll}
\toprule
\textbf{Module} & \textbf{Used in} & \textbf{Role} \\
\midrule
\code{tkinter}, \code{tkinter.ttk} & \file{gui.py} & the entire interface \\
\code{math} & \file{morse\_audio.py} & \code{sin} and \code{pi} for the sine wave \\
\code{struct} & \file{morse\_audio.py} & pack sample numbers into raw bytes \\
\code{wave} & \file{morse\_audio.py} & write the WAV header and frames \\
\code{shutil} & \file{morse\_audio.py} & \code{which} to find a player; \code{rmtree} to clean up \\
\code{subprocess} & \file{morse\_audio.py} & launch the player without blocking \\
\code{tempfile} & \file{morse\_audio.py} & a scratch directory for the two tones \\
\code{atexit} & \file{morse\_audio.py} & delete that directory on exit \\
\code{os} & \file{morse\_audio.py} & join paths portably \\
\code{sys} & \file{morse\_audio.py} & \code{sys.platform} to detect Windows \\
\bottomrule
\end{tabular}
\caption{Every dependency, and why it is there. \file{morse\_converter.py} imports
nothing at all.}
\end{table}

There is no configuration file, no environment variable, no network access, and no
persistence: the program writes only two temporary WAV files, and deletes them on
exit. The tuning knobs (\code{UNIT\_MS}, \code{TONE\_FREQUENCY\_HZ}, \code{VOLUME},
\code{SAMPLE\_RATE}, and the colour constants) are module-level constants, editable
only by editing the source. For a project this size that is a reasonable stopping
point, not a gap.

\file{.gitignore} excludes the usual Python build and cache artefacts
(\code{\_\_pycache\_\_/}, \code{*.pyc}, virtual environments, coverage output) plus
\code{.env}, ensuring a credentials file could never be committed by accident. The
licence is PolyForm Strict 1.0.0: source-available for viewing, with no permission
to use it in a product, which is the appropriate choice for a portfolio piece.

% =====================================================================
\section{What I verified by running it}

Nothing in this document about behaviour is taken on trust from the README. The
following were executed against the code in this folder:

\begin{itemize}
\item \textbf{Table integrity.} \code{MORSE\_CODE} has 57 entries;
      \code{TEXT\_CODE} has 57; zero Morse patterns are duplicated, so the inversion
      loses nothing.
\item \textbf{Round trip.} \code{morse\_to\_text(text\_to\_morse(s)) == s.upper()}
      held for every input tested, including mixed case, punctuation, digits, a
      literal \code{/} in the text, leading/trailing spaces, doubled internal spaces,
      and an all-whitespace string.
\item \textbf{The redundancy claim.} A naive one-line decoder that never splits on
      \code{" / "} produced identical output to \code{morse\_to\_text} on every test
      input, confirming the outer split adds no behaviour.
\item \textbf{Timing.} \code{build\_timing\_sequence} for \code{"SOS"}
      (\code{... --- ...}) emits 17 events totalling 27 units, and for
      \code{"SOS SOS"} 35 events totalling 61 units. \code{render\_morse\_to\_wav}
      produced files of 95{,}256 and 215{,}208 frames at 44{,}100~Hz, exactly
      2160~ms and 4880~ms, matching $27 \times 80$~ms and $61 \times 80$~ms computed
      by hand from the standard ratios. The timing claim is true.
\item \textbf{Backend detection.} \code{detect\_playback\_backend()} returned
      \code{'aplay'} on this machine, matching the README.
\item \textbf{The newline limitation.} \code{text\_to\_morse("A\textbackslash n")}
      raises \code{ValueError}, as does a tab, as does an accented letter.
\item \textbf{Silent garbage tolerance.} \code{build\_timing\_sequence(".- XYZ
      -...")} returns a timeline with \code{XYZ} dropped and no error raised.
\end{itemize}

I did not run the GUI itself (this session has no display), so statements about
on-screen behaviour are read from the code rather than observed. The project
includes a screenshot at \file{docs/screenshots/app.png} showing the running window.
I also could not test the \code{winsound}, \code{afplay} or \code{paplay} paths,
which need Windows, macOS, or PulseAudio respectively; the README already states
this limitation itself.

% =====================================================================
\section{Honest summary}

What is genuinely good here:

\begin{itemize}
\item The separation of logic, audio and interface is real, not nominal, and it is
      the reason the audio module could be rewritten without touching the converter.
\item The single event list feeding both the visual and the audio timeline is the
      right answer to a problem the project got wrong the first time, and the README
      documents that history instead of hiding it.
\item Runtime backend detection with a genuine no-audio fallback is more care than a
      practice project usually gets.
\item Generating real sound from \code{math.sin}, \code{struct} and \code{wave} is a
      neat, correct, and well-earned way to keep the dependency list empty.
\item The verified timing holds to the millisecond.
\end{itemize}

What is weak, in order of how much it matters:

\begin{enumerate}
\item \textbf{No automated tests.} For a project whose correctness lives in a
      57-entry table, the absence of a test file is the biggest gap. Every claim in
      this document had to be re-established by hand. The README names this itself as
      the first thing it would do differently, which is the right instinct.
\item \textbf{\code{\_current\_tab\_state} is misnamed and ignores the selected
      tab}, producing genuinely surprising Play behaviour when both tabs have content.
\item \textbf{\code{build\_timing\_sequence} silently drops invalid symbols} while
      \code{morse\_to\_text} raises on them: two different contracts for the same
      malformed input.
\item \textbf{The redundant \code{" / "} split}, defended in the README by an
      argument that does not address the simpler alternative.
\item \textbf{\code{\_safe\_text\_to\_morse} and \code{\_safe\_morse\_to\_text} are
      empty wrappers} whose names promise safety they do not deliver.
\item \textbf{Multi-line input boxes that reject newlines}, handled gracefully but
      still a mismatch between the widget and the alphabet.
\item \textbf{\code{render\_morse\_to\_wav} is unused by the application}, and the
      unreaped \code{Popen} processes leave short-lived zombies.
\end{enumerate}

None of these is a crash, a security issue, or a wrong answer. They are the ordinary
residue of a project that grew from a script into an application, and the README is
unusually candid about most of them already. The gap between the code and its
documentation is small, and where it exists (the \code{" / "} rationale) it is a
matter of a defence that no longer fits rather than a false claim about what the
software does.

\end{document}
